\chapter{Conclusion}

This thesis investigates skill representation for scalable agent skill retrieval. The motivating problem is semantic confusion: as skill libraries grow, many skills can sound relevant while requiring different procedures. In such cases, the correct skill may exist, but the retrieval system may surface a plausible yet procedurally inappropriate alternative.

The thesis reframes skill selection as a layered problem involving skill artifacts, selection representations, and retrieval policies. This framing is important because a skill artifact may contain rich procedural information while the selector sees only a compressed card. Conversely, a strong retriever may still fail if the representation hides the distinctions needed for procedural selection.

The current benchmark contains 2401 skills and 137 controlled prompts, with generated scale skills, imported public skills, and a manually adjudicated public-gold stratum. Current evidence suggests that flat metadata and generic full-skill embedding alone are fragile under scale. SkillRouter provides a strong skill-specific retrieval baseline: its embedding plus reranker reaches 83.2\% top-1 on the controlled benchmark. The strongest current controlled result combines SkillRouter first-stage retrieval with M6-v1-local task/output/workflow reranking, reaching 86.9\% top-1, 98.5\% top-5, and 0.927 MRR. However, public-gold results show that naive schema overlap and the first deterministic M6-v1-local matcher are not enough when skills are externally authored, broad, duplicated, or implicit.

The main interim answer to RQ1 is that skill representations should preserve at least use conditions, expected outputs, workflow/procedure, inputs or preconditions, boundaries, and conditional dependency/resource information. The main interim answer to RQ2 is that dense retrieval is useful for broad candidate generation, especially when using a skill-specific model such as SkillRouter, but final selection benefits from procedural reranking when the procedural information is accurately extracted and used field-to-field.

The next technical step is to improve M6-v1 beyond the deterministic local prototype by adding semantic or LLM-assisted field matching, especially for public-authored skills. If time permits, M4 tree routing and M5 graph retrieval should be evaluated as structure-aware architecture comparisons. The final thesis should then add cost/latency analysis, cleaned public-gold results, and a small downstream validation study.

The expected contribution is not a universal skill-router benchmark or a provider leaderboard. It is a representation-centred account of why semantically similar skills are difficult to retrieve, which procedural information appears to matter, and how retrieval architectures can use that information to make large agent skill libraries more reliable.
